% ============================================================================
% TEMP DRAFT (Claude + Codex synthesis) — condensed Intro + Results opening.
% Included before the real Introduction for review. Delete this file and its
% \input in main.tex once a decision is made.
% ============================================================================

\section{[DRAFT] Introduction (condensed)}

Clinical and lesion studies suggest that declarative memories initially depend on the
hippocampus, but are later transferred to other brain areas \cite{Dudai2015, Squire2015,
Sekeres2017}. Some forms of memory eventually become independent of the hippocampus and
depend only on a stable representation in the neocortex \cite{Dudai2015, Squire2015,
Sekeres2017}. Similarly, procedural memories are consolidated within cortico-striatal
networks \cite{Karni1994, Brashers-Krug1996, Dudai2015}. This process of memory
transformation --- termed systems memory consolidation --- is thought to prevent newly
acquired memories from overwriting old ones, thereby extending memory retention times
(``plasticity-stability dilemma''; \cite{Grossberg1987, Abraham2005, Fusi2005,
Leibold2008, Roxin2013}), and to enable a simultaneous acquisition of episodic memories
and semantic knowledge of the world \cite{McClelland1995, Kumaran2016}.

While specific neuronal activity patterns, including accelerated replay of recent
experiences \cite{Lee2002, Skaggs1996}, are involved in the transfer of memories from
hippocampus to neocortex \cite{Diekelmann2010}, the mechanisms underlying systems memory
consolidation are not well understood. Specifically, it remains unclear how this
consolidation-related transfer is shaped by the anatomical structure and the plasticity of
the underlying neural circuits. This poses a substantial obstacle for understanding into
which regions memories are consolidated; why some memories are consolidated more rapidly
than others \cite{Nadel1997, Tse2007, Brodt2018}; why some memories stay hippocampus
dependent, and why and how the character of memories changes over time \cite{Dudai2015};
and whether the consolidation of declarative and non-declarative memories \cite{Karni1994,
Brashers-Krug1996, Dudai2015} are two sides of the same coin. These questions are hard to
approach within phenomenological theories of systems consolidation such as the standard
consolidation theory \cite{Squire1995, McClelland1995}, the multiple trace theory
\cite{Nadel1997}, and the trace transformation theory \cite{Winocur2010, Winocur2011}.

Here, we propose a novel mechanistic foundation of the consolidation process. Our focus
lies on simple forms of memory that can be phrased as feedforward cue-response
associations, stored in synaptic pathways between an input area --- neurally representing
the cue --- and an output area --- neurally representing the response. Our work thus
relates to hetero-associative memory, and is therefore applicable to both declarative and
non-declarative memories, rather than to recurrent, auto-associative memory (see, e.g.,
\cite{Hopfield1982, Zenke2015, Tome2020}). Our central hypothesis --- the parallel pathway
theory (PPT) --- is that systems memory consolidation arises naturally from the interplay
of two abundantly found neuronal features: parallel synaptic pathways and Hebbian
plasticity \cite{VanEssen1992, Malenka2004}. We illustrate this theory in a simple
hippocampal circuit motif, analyse its basic conditions, and show in simulations that it
is robust to neuronal nonlinearities, accounts for a hippocampal lesion study in rodents
\cite{Remondes2004}, and --- iterated in a cascade --- achieves full consolidation into
neocortex with power-law forgetting \cite{Wixted2004} and a semantisation of episodic
memories.

\section{[DRAFT] Results opening (condensed)}

\subsection{A mechanistic basis for systems memory consolidation}

The parallel pathway theory (PPT) conceptualises a memory as a cue-response association
stored in a synaptic pathway between an input area and an output area. The theory relies
on a common circuit motif in which a direct, monosynaptic pathway and an indirect,
multisynaptic pathway converge onto the same output area. Memories are initially stored in
the indirect pathway and are subsequently transferred to the direct pathway by Hebbian
plasticity. Because the indirect pathway transmits activity with a longer delay, inputs in
the direct pathway precede the output spikes driven by the indirect pathway. A
timing-dependent plasticity rule can therefore use the indirect pathway as a teacher for
the direct pathway (Fig~\ref{fig:PPT_mechanism}A).

In the hippocampal formation, this motif maps onto two EC inputs to CA1: the direct,
monosynaptic perforant path (\ppca, Fig~\ref{fig:PPT_mechanism}B, red) and the indirect,
trisynaptic pathway through DG and CA3, reaching CA1 through the Schaffer collaterals
(SC; Fig~\ref{fig:PPT_mechanism}B, blue; \cite{Amaral1993}). As in earlier theories, we
assume that the indirect pathway via CA3 is involved in the original storage of memories
\cite{Marr1969, Treves1994}, consistent with experimental evidence \cite{Brun2002,
Nakazawa2003, Nakashiba2008}. A specific neural activity pattern in EC --- the cue --- can
then trigger spikes in a subset of CA1 cells through the SC, representing the associated
response. The same cue also reaches CA1 through the \ppca, but this direct input initially
fails to trigger spikes because its synaptic weight pattern does not yet match the cue.
However, cue-driven \ppca~inputs precede the CA1 spikes triggered by the indirect pathway
by 5--15~ms due to transmission delays \cite{Yeckel1990}. Presynaptic spikes preceding
postsynaptic spikes with such a short delay favour selective long-term potentiation by
spike timing-dependent plasticity (STDP, Fig~\ref{fig:PPT_mechanism}C) \cite{Bi1998,
Markram1997, Gerstner1996}. Consequently, cue-driven \ppca~synapses onto activated CA1
cells are strengthened until the memory that was initially stored in the indirect pathway
can be recalled via the direct pathway alone.

To illustrate this mechanism, we used a simple integrate-and-fire neuron model
(Methods Sec.~\ref{ssec:ppt_methods_LIF}) of a CA1 cell receiving inputs through the SC
and the \ppca. The two pathways contained the same number of synapses and transmitted
identical spike patterns apart from a 5-ms delay in the SC (Fig~\ref{fig:PPT_mechanism}D),
so that consolidation corresponded to copying the synaptic weight pattern of the SC to the
\ppca. In line with our hypothesis, STDP in the \ppca~synapses achieved this consolidation
(Fig~\ref{fig:PPT_mechanism}E), whereas consolidation in the opposite direction failed
because the temporal order of spiking activity is reversed and hence does not favour
synaptic potentiation (Fig~\ref{fig:PPT_mechanism}F).

To understand the conditions under which the PPT can consolidate associative memories, we
performed a mathematical analysis of the mechanism. It makes two key predictions: STDP in
the direct pathway implements a linear regression that approximates the input-output
mapping of the indirect pathway (Sec.~\ref{ssec:theory_stdp} and
\ref{ssec:theory_dynamics}), so the mechanism should generalise to differing cue
representations in the two pathways; and consolidation is most effective when the
correlation time of the input is matched to the coincidence time scale of STDP
(Sec.~\ref{ssec:theory_timing}, Fig~\ref{sfig:theory_timing}B). In the following sections,
we test these predictions in simulations, show that the mechanism is robust to neuronal
nonlinearities, that it accounts for lesion studies in rodents, and that it can be
hierarchically iterated.
